\section{Structural search: changing the statement that is being proved}
\label{sec:structural}

This is the capability that most clearly exceeds local saturation. A saturation
engine is very effective when useful ground terms and useful lemmas are already
present. A much wider class of goals requires \emph{making one of those things
exist}.

\subsection{Demand-directed theorem use}

Trigger-based instantiation works forward, from terms present in the equality
state. The documented E-matching machinery is a baseline, not something to
discard~\cite{grind-ematch}. \forge{} adds a complementary backward direction.

Maintain an index over the \emph{conclusions} of approved theorems --- initially
a discrimination tree keyed by the target's head symbol and normalised type
structure. An entry records the theorem constant, its universe and type
parameters, its conclusion head, a discrimination key, the shapes of its
premises, and a cheap estimate of elaboration cost. Retrieve a bounded candidate
list, instantiate universes, and unify the theorem conclusion with the target.
Successful unification creates premises as AND obligations and instantiation
choices as OR alternatives, with explicit dependencies so that the conclusion
cannot become available before its premises are proved.

For a theorem of shape $\forall x,\;P(x)\to Q(f(x))$ and a goal $Q(f(u))$,
backward application immediately proposes the obligation $P(u)$ and introduces
useful terms even when no suitable trigger was present in the initial forward
state. For a goal $Q(t)$ with arbitrary $t$ it does \emph{not} justify inventing
an inverse of $f$: the system must unify $t$ with an application already known
to be equal, synthesise a candidate $u$ and prove $f(u)=t$, or decline.

\begin{remark}
This is not a claim that E-matching could never reproduce the same argument with
different patterns and annotations. What a conclusion index removes is the need
for a human to have chosen those patterns in advance. Conclusion-selected
backward patterns already exist in \grind{} in some form~\cite{grind-paper}; the
proposed addition is that a retrieved theorem's unproved premises become
persistent obligations that other workers can be scheduled against.
\end{remark}

Premise retrieval should combine the target with currently \emph{blocked}
obligations, not just a bag of symbols from the original goal. A theorem
resolving a denominator guard can be more useful than one mentioning the final
function symbol. Start with local hypotheses, the equations of mentioned
definitions, relevant extensionality lemmas, and a small retrieved set; expand
only if those fail.

\subsubsection*{Solving for the missing term}

Backward matching does not determine every parameter, and the honest responses
to that are limited. In $\forall x\,y,\;R(x,y)\to P(x)$, a demand for $P(a)$
leaves $y$ unknown. \forge{} must not silently choose an arbitrary inhabitant.
It can enumerate appropriately typed terms from the local context, ask a witness
engine, exploit equalities, or leave the application suspended. Every selected
instantiation is checked by Lean, and typeclass synthesis and inhabitation
obligations are scheduled explicitly rather than assuming an arbitrary type has
an element.

A sharper technique applies when matching stalls at an \emph{interpreted}
subterm. Rather than failing, emit an equality constraint between that subterm
and the corresponding source expression, and solve for the missing binder.

\begin{lstlisting}
matchDemand(theorem, goal):
  freshen theorem parameters in a child metavariable context
  unify rigid heads and directly matching arguments
  collect the remaining argument equations
  dispatch recognized affine equations to witness synthesis
  enumerate small typed terms only for the still-free parameters
  apply the theorem with a complete candidate assignment
  retain every generated side condition as an obligation
  accept only after all obligations and the final proof replay
\end{lstlisting}

For hypotheses $\forall k\in\Z,\;P(2k+1)$ and $n\bmod 2=1$ with goal $P(n)$, the
controller synthesises $k=n/2$ from the arithmetic shape and proves
$2(n/2)+1=n$ using the oddness assumption, then transports along that checked
equality. For a conclusion containing $f(t+1)$ against a goal involving $f(s)$,
the candidate $t=s-1$ comes from solving $t+1=s$ over $\Z$ --- and over $\N$ the
same surface subtraction has different semantics and requires a guard.

Two disciplines make this safe. An arithmetic identity used to match arguments
must appear in the proof as rewriting or transport; it cannot be smuggled into
the unifier as an unverified definitional equality. And several triggers or
repeated binders produce a \emph{joint} constraint system: solving each
occurrence independently loses the equality relating two occurrences of the
same variable. Cache instances by theorem identity, universe and type
instantiation, the complete substitution, and the valid branch context. A proof
of $P(x)$ under $h:x=0$ is not a proof of the same displayed goal under
$h:x=1$.

\subsection{Demand slicing in a finite Horn fragment}
\status{Implemented and tested in Python by four independent efforts, in both
propositional and first-order ground variants.}

The prototypes isolate the mechanism in a fragment small enough to prove things
about: finite ground definite Horn programs. Backward relevance closure selects
the rules that could contribute to the target; forward chaining then runs only
on that slice.

\begin{proposition}[Slicing preserves derivability]
\label{prop:slicing}
For a finite ground definite Horn program, restricting forward chaining to the
backward conclusion-dependency closure of a target preserves whether that target
belongs to the least model.
\end{proposition}
\begin{proof}
A derivation of the target is a finite tree of rule applications. Starting from
its root, every rule used and every premise of every such rule is included by
the backward closure, so every derivation of the target is retained. Conversely
every rule of the restricted program belongs to the original program, so any
restricted derivation is valid in the original. Forward fixed-point evaluation
computes derivability in both finite programs. Cycles without a supporting fact
create no derivation in either.
\end{proof}

Termination and non-circularity are handled by \emph{tabling}, not by a
depth-first convention that treats ``currently being visited'' as success. Each
rule instance carries a counter of unresolved premises; input facts initialise
the agenda; proving a premise decrements the counters of its dependents; a rule
fires only when every premise is proved, and the proof node is stored at that
moment. Rules $P\to Q$ and $Q\to P$ with no base fact prove neither.

\begin{proposition}[Relative completeness of the finite core]
For a fixed finite set of ground candidate rule instances, if the dependency
agenda is exhausted without a resource cutoff, every fact in their least forward
closure is eventually proved. In particular, a backward-relevant closure
containing a ground derivation of the target will discover a proof.
\end{proposition}
\begin{proof}
Induct on the height of a finite derivation. Height-zero input facts enter the
agenda initially. Once all smaller-height premises of an instance have entered,
its unresolved counter reaches zero and its conclusion is enqueued. Tabling
ensures that repeated discoveries do not obstruct propagation.
\end{proof}

The certificate is a topologically ordered DAG: every premise index is strictly
smaller than its conclusion's index; the substitution domain must be exactly the
head's variables; and the head instantiation and premise alignment are checked.
A dependency cycle in the \emph{search} graph may be legitimate, but the final
proof certificate cannot justify itself through that cycle.

Two honesty constraints govern how this is measured, and
Section~\ref{sec:experiments} keeps four separate Horn experiments apart for
exactly this reason. First, the cost of \emph{building} the conclusion index
over all input rules is part of the demand-directed method's cost; none of the
prototypes obtains a sublinear end-to-end runtime merely by activating fewer
rules. Second, proof minimisation reduces even a full-closure derivation to the
same small proof, so \emph{proof size alone conceals the search-work
difference} --- the interesting quantity is rule firings, not the length of the
answer.

\begin{remark}
None of these Horn engines is an implementation or emulation of \grind{}'s
E-matching. They compare two deliberately small search policies against each
other, and one of them is a fair semi-naive forward-chaining baseline rather
than a straw man.
\end{remark}

\subsection{Induction: selecting a scheme and generalising the statement}
\status{The list and natural-number fragments, changing-argument detection,
bounded right-hand-side synthesis, structural induction, and trace replay are
implemented. General dependent induction and e-graph-guided frontier discovery
are specified only.}

The induction planner starts from the recursive applications visible in the
target and in relevant hypotheses. For each function it records which argument
positions are destructed by pattern matching, which are passed to recursive
calls as proper subterms, and which non-structural arguments change. Candidates
are ranked by how many blocked recursive calls they expose after one constructor
split. A variable is not selected merely because its type has an induction
principle, and blind induction on every natural variable is not acceptable.

A case split on an inductive value does not automatically supply an induction
hypothesis; confusing the two is a design error. Each candidate plan carries an
explicit Lean motive, an induction principle, the generalised variables, and the
recursive hypotheses available in each case. User-specified schemes take
priority.

Generalisation follows a \emph{dependency closure}, and this is where a naive
implementation goes wrong. In a context containing $x:A$, $y:B(x)$, and
$h:C(x,y)$, induction on $x$ requires a motive of the form
\[
  M(x)\;:=\;\forall y:B(x),\; C(x,y)\to P(x,y),
\]
not a proposition containing the old fixed $y$ at every new value of $x$. If a
generalised variable occurs in the type of another binder, that binder must be
generalised or explicitly transported as well. A textual substitution that
ignores dependent hypotheses is not a correct generalisation algorithm, and a
well-typed-looking first-order erasure is not a substitute for building the
motive through Lean's elaborator.

Generalisation is an \emph{OR alternative}, not an irreversible preprocessing
step: too much generality can make a useful theorem false or unnecessarily
difficult. Candidate generalisation sets are ordered by size and capped.

\subsubsection*{The accumulator example, including the obstruction}

Write $A$ for append, $R$ for reverse, and $\mathrm{RA}$ for accumulator
reverse:
\begin{align*}
A([],y)&=y,& A(h::t,y)&=h::A(t,y),\\
R([])&=[],& R(h::t)&=A(R(t),[h]),\\
\mathrm{RA}([],a)&=a,& \mathrm{RA}(h::t,a)&=\mathrm{RA}(t,h::a).
\end{align*}
The requested theorem is $\mathrm{RA}(x,[])=R(x)$. Inducting on $x$ without
generalisation gives an induction hypothesis about $\mathrm{RA}(t,[])$, whereas
the recursive call in the step case is $\mathrm{RA}(t,[h])$. \emph{Increasing a
rewriting budget cannot make this particular hypothesis say something
stronger.} This is the shape of obstruction the planner exists to detect.

The changing-argument detector observes that argument~0 is destructed
structurally while argument~1 changes from $a$ to $h::a$. It replaces the fixed
empty accumulator with a fresh parameter and asks for a right-hand side in a
typed grammar
\[
  L ::= []\;\mid\; x\;\mid\; a\;\mid\; R(L)\;\mid\; A(L,L),
\]
enumerated in increasing syntax-tree size. The successful candidate is
\begin{equation}\label{eq:accspec}
\forall x\,a,\qquad \mathrm{RA}(x,a)=A(R(x),a).
\end{equation}
The base case is immediate. The step case uses the generalised hypothesis
$\mathrm{RA}(t,h::a)=A(R(t),h::a)$ against
\[
A(R(h::t),a)=A(A(R(t),[h]),a)=A(R(t),A([h],a))=A(R(t),h::a),
\]
the middle step being append associativity. Instantiating \eqref{eq:accspec} at
$a=[]$ and using append's right identity proves the original theorem.

A second family makes the same point over $\N$. For
\[
 A(0,a)=a,\qquad A(n+1,a)=A(n,a+2n+1),
\]
the theorem $A(n,a)=a+n^2$ needs $a$ generalised before induction; the step then
reduces to $a+2n+1+n^2=a+(n+1)^2$. The automation problem is choosing the
generalisation, not performing the final ring calculation.

\begin{remark}[Sampling rejects; it never proves]
Two prototypes filter grammar candidates against finite sample environments ---
one uses all 49 pairs of lists of length at most two over a two-element
alphabet. Samples \emph{only reject}. The surviving candidate is accepted solely
because it then passes symbolic induction and independent trace replay.
\end{remark}

\subsubsection*{What a replay checker must restrict}

The miniature checkers make the induction discipline a first-class, checked
datum, and this is the part a production implementation most needs to inherit.

There are separate syntactic constructors for universal pattern variables and
for \emph{rigid} free variables. In an induction step the recursion index --- a
list tail, a natural predecessor --- is rigid; only genuinely generalised
parameters remain instantiable in the induction hypothesis. Treating the tail as
an ordinary rewrite metavariable would let an induction hypothesis rewrite the
goal at an unjustified larger argument. In the base case the induction
hypothesis is not available at all.

The strongest formulation attaches to each rule an explicit \emph{flexible}
variable set --- the variables at which that rule, notably the induction
hypothesis, may be instantiated --- and requires each replayed step to supply
exactly that substitution set with correct sorts, rejecting missing or illegal
instantiations.

Every rewrite certificate contains a rule identifier and a subterm path. The
checker retrieves the rule from the defining equations, from an earlier checked
lemma, or from the current induction hypothesis; matches the left-hand side;
checks repeated-variable identity and sorts; replaces the selected subterm; and
continues. The two replayed sides must end syntactically identical. No search
and no normalisation happen during replay.

\begin{proposition}[Soundness of the list certificate format]
Assume that previously installed equations are valid for all finite lists. If
the checker accepts an induction certificate, its universally quantified
equation holds for all finite lists under the supplied definitions.
\end{proposition}
\begin{proof}
Each defining rewrite preserves the interpreted value. Each installed lemma
instance preserves it by assumption, and congruence permits replacement inside
a typed term context. An accepted base trace therefore proves the empty-list
case. In the step trace the only additional rule is the desired statement for
the rigid tail, universally quantified over its generalised parameters, and that
rule is valid under the structural induction hypothesis. The accepted step trace
therefore proves the constructor case, and structural induction establishes the
universal equation.
\end{proof}

This justifies a small format. It is not a machine-checked Lean proof of the
Python implementation, and the negative tests reduce rather than remove the risk
of a checker bug.

\subsubsection*{Lemma discovery and non-circularity}

A public theorem-installation operation first checks both induction cases, and
\emph{a lemma is not available while its own certificate is being checked} ---
which rejects a forged proof that cites its own conclusion. A bundle of lemmas
is accepted only in a dependency order in which no entry refers to a later one,
and shadowing a defining equation or the induction-hypothesis name is rejected
outright.

Beyond generalisation, a demand-driven lemma explorer works from the residual
of a failed step. Retain the normalised left and right frontiers and the
rewrites that produced them. Bounded equality saturation can expose alternative
forms a one-way normaliser would have discarded --- for the accumulator example,
unfolding append in reverse exposes $h::a$ as $A([h],a)$, and the remaining
frontier then suggests $A(A(u,v),w)=A(u,A(v,w))$.

\begin{lstlisting}
find_lemma(left_frontier, right_frontier):
  expand each frontier using a bounded set of proved equalities
  identify pairs separated by a small structural difference
  anti-unify shared subterms, preserving repeated-variable identity
  abstract a well-typed parameter telescope
  reject trivial, cyclic, or irrelevant candidates
  filter candidates with executable counterexamples where available
  prove surviving candidates in a separate obligation
  return only checked lemmas to the main equality state
\end{lstlisting}

Typed anti-unification starts in a restricted fragment: fully applied,
homogeneous first-order terms with compatible types. Memoise pairs of subterms
so that anti-unifying $f(a,a)$ and $f(b,b)$ gives $f(z,z)$ and not $f(z_1,z_2)$.
In a dependent fragment, an abstraction that changes a term's type needs a typed
telescope and possibly transport proofs; failing to find one is an ordinary
search failure.

Three further candidate sources are worth implementing: generalise a changing
non-structural argument; propose a distribution or homomorphism lemma such as
$R(A(xs,ys))=A(R(ys),R(xs))$ when nested calls obstruct normalisation; and
synthesise a stronger invariant containing the measure, bound, or relation the
step requires. Useful templates for the last are
\[
 f(s,a)=H(s,a),\qquad
 \mu(f(s,a))=\mu(a)+J(s),\qquad
 R(f(s,a),s,a).
\]
Each candidate becomes its own obligation, never a rewrite rule merely because
it helped simplify another candidate. Quotient syntactically different
candidates only by \emph{already proved} equalities; otherwise the quotient can
discard the correct expression or conceal an unproved assumption.

This goal-directed exploration is closely related to CCLemma's use of e-graphs
to focus inductive lemma discovery on useful proof frontiers, and to the
rippling tradition; the proposal adapts that work to a proof-producing
controller rather than originating
it~\cite{cclemma,rippling,directed-lemma-synthesis}.

Impose separate limits on term size, saturation nodes, candidate count,
induction depth, and lemma-dependency depth, and use term fingerprints and
alpha-normalised telescopes to avoid reproving the same lemma under new binder
names. Bounded enumeration is exponential in its grammar size; the production
advantage must come from recursive-call analysis, relevance filtering, typed
indexing, and frontier selection, not from assuming that unrestricted
enumeration will be cheap.

\subsection{Invariant synthesis by exact CEGIS}
\status{Implemented and tested in Python by six independent efforts, in
univariate, bivariate, multivariate, accumulator and state-machine forms.}

\subsubsection*{The additive fragment and its certificate}

Consider
\begin{equation}\label{eq:recurrence}
T(0)=t_0,\qquad T(n+1)=T(n)+f(n),\qquad n\in\N,
\end{equation}
with $t_0\in\Q$ and $f\in\Q[X]$. Search for $q\in\Q[X]$ with
\begin{equation}\label{eq:rec-cert}
q(0)=t_0,\qquad q(X+1)-q(X)=f(X).
\end{equation}
The certificate is the exact coefficient list of $q$. Its checker computes the
shift $q(X+1)$ by polynomial arithmetic and compares exact coefficients; it does
not sample the step relation and does not trust an interpolator.

\begin{proposition}[Recurrence certificate]
\label{prop:rec}
If \eqref{eq:rec-cert} holds then $T(n)=q(n)$ for every $n\in\N$.
\end{proposition}
\begin{proof}
The base equality is the first condition. If $T(n)=q(n)$ then
\eqref{eq:recurrence} and the second condition give
$T(n+1)=q(n)+f(n)=q(n+1)$. Apply induction on $n$.
\end{proof}

The multivariate form carries parameters. For
$S(0,\bm a)=I(\bm a)$ and $S(n+1,\bm a)=S(n,\bm a)+P(n,\bm a)$, search for
$F(n,\bm a)=\sum_{|\alpha|\le D}c_\alpha(n,\bm a)^\alpha$, substitute into
$F(0,\bm a)=I(\bm a)$ and $F(n+1,\bm a)-F(n,\bm a)=P(n,\bm a)$, and equate
coefficients to obtain a rational linear system in the $c_\alpha$.

For the accumulator shape
$\mathrm{go}(0,a)=a$, $\mathrm{go}(n+1,a)=\mathrm{go}(n,a+r(n))$, the ansatz is
$F(n,a)=P(n)+aQ(n)$ with obligations $F(0,a)=a$ and $F(n+1,a)=F(n,a+r(n))$.
One implementation joins the two systems by introducing a formal tag variable
$t$, encoding the base equation in the coefficients of $t^0$ and the step
equation in those of $t^1$, so that a single exact row reduction solves both
simultaneously; $t$ does not appear in the final invariant.

\subsubsection*{Counterexamples direct the coefficient search}

For degree bound $d$, write $q(X)=\sum_{i=0}^d c_iX^i$, maintain exact sample
constraints $q(n)=T(n)$, and solve the rational linear system with free
coefficients set to zero. Check the candidate's base and step identities. If the
step residual
\[
r(X)=q(X+1)-q(X)-f(X)
\]
is nonzero, choose a natural $n$ with $r(n)\ne0$ and add \emph{both} $q(n)=T(n)$
and $q(n+1)=T(n+1)$ to the sample set. At least one is a new constraint: if both
were already satisfied by the candidate, their difference would force $r(n)=0$.
A nonzero degree-$k$ polynomial cannot vanish at all of $0,1,\ldots,k$, so an
exact separating counterexample is found without an SMT solver. If the
coefficient constraints become inconsistent, raise the degree bound.

\begin{proposition}[Fragment completeness]
\label{prop:rec-complete}
For a nonzero $f$ of degree $k$, the finite-difference map
$q\mapsto q(X+1)-q(X)$ sends a leading term $aX^{k+1}$ to $(k+1)aX^k$ plus lower
terms. Over $\Q$ one may therefore solve for the leading coefficient, subtract
its difference, and continue downward: a degree-$(k+1)$ solution exists and is
unique once $q(0)$ is fixed. The zero increment gives a constant solution.
Hence exact linear solving with a degree budget of at least $k+1$ terminates on
this fragment.
\end{proposition}

Equivalently, the forward-difference operator maps polynomials of degree at most
$d+1$ with zero constant term bijectively onto polynomials of degree at most
$d$, since its matrix in the monomial basis is triangular with nonzero diagonal.
In the binomial basis $\Delta\binom{n}{k+1}=\binom{n}{k}$, which a future
implementation could exploit; the prototypes use ordinary monomials and exact
row reduction because that keeps the certificate checker small.

This is a \emph{fragment-level} coverage result. It says nothing about
recurrences involving products with $S$, branching updates, nonlinear state
transformations, piecewise definitions, or integer formulas with truncated
division, and it does not make the broader induction planner complete.

\subsubsection*{The overfitting trap}

For $f(X)=X+1$ the discovered result is $q(X)=X(X+1)/2$. The interpolated
expression
\[
\widetilde q(X)=\frac{X(X+1)}2+X(X-1)(X-2)(X-3)
\]
agrees with the triangular-number formula on the first four inputs and is
rejected by the step checker. This is a small, concrete reason why conjecture
generation and proof acceptance must be separate operations, and the same trap
is built into the test suites as a hostile control: a high-degree term that
vanishes at every sampled point but violates the recurrence identity elsewhere.

For $P(n)=n^2$ with $I=0$ the synthesised result is
\[
 F(n)=\frac{n^3}{3}-\frac{n^2}{2}+\frac n6=\frac{n(n-1)(2n-1)}6,
\]
and one run recovered the degree-9 invariant for $r(n)=n^8$,
\[
F(n,a)=a+\frac{n^9}{9}-\frac{n^8}{2}+\frac{2n^7}{3}
 -\frac{7n^5}{15}+\frac{2n^3}{9}-\frac{n}{30}.
\]

\subsubsection*{The state-machine generalisation}

A more reusable certificate format treats the invariant as a conserved quantity
of a transition system: for initial state $s_0$ and transition $T$, require
\[
 I(s_0)=0,\qquad I(T(s))=I(s).
\]
The second identity is a convenient \emph{sufficient} condition --- proving
$I(s)=0\Rightarrow I(T(s))=0$ would also suffice but is not the implemented
certificate language. For sums of squares, clearing denominators gives
$J(n,s)=6s-n(n+1)(2n+1)$, and
\[
 J(n+1,s+(n+1)^2)-J(n,s)
 =6(n+1)^2-(n+1)(n+2)(2n+3)+n(n+1)(2n+1)=0 .
\]

The next step is a genuine invariant worker for a state transition
$s'=F(n,s)$, proving the initial condition and
\[
 I(n,s)\land\mathrm{Guard}(n,s)\;\Longrightarrow\;I(n+1,F(n,s)).
\]
Coefficient comparison handles identities and exact nonlinear certificates can
handle some inequalities. A warning applies to the multi-invariant version: if
the step condition takes the form $I_i(T(\vec x))=\sum_j h_{ij}(\vec x)I_j(\vec
x)$, then jointly discovering the $I_j$ and the multipliers is a
\emph{bilinear} search, not linear algebra, and it should not be described as
the latter. Separate candidate generation from multiplier certification, or fix
one side at a time.

\begin{remark}
An invariant does not prove that a loop or a recursive definition terminates.
Termination is a separate obligation.
\end{remark}

\subsection{Witness synthesis: construct terms, not satisfying assignments}

\subsubsection*{The logical distinction}

To prove $\forall x\,\exists y\,P(x,y)$, a satisfying assignment for one value
of $x$ is not enough. A useful synthesiser produces a term $y=w(x)$ and proves
$\forall x\,P(x,w(x))$. For
\[
\forall x\in\R,\quad \exists y\in\R,\quad x+1\le y\le x+2,
\]
the witness is $w(x)=x+1$; no constant works for every real $x$, and a
constant-only template returns no certificate where the affine template
succeeds.

Sampling individual values of $x$ and finding unrelated values of $y$
constructs no function and proves nothing. The grammar must represent a
well-typed expression $w(x)$. A synthesis oracle may help choose $w$; it cannot
justify $P(x,w(x))$ by testing.

Witness binders retain their dependencies. A witness for the second existential
may refer to an earlier witness but not to a future binder or to a branch-local
variable outside its scope. The instructive negative example is
$\exists x:\alpha,\forall y:\alpha,\;x=y$: the witness for $x$ cannot be the
later-bound $y$, and a naive synthesiser that sees all displayed variables will
propose exactly that. The existential task's context must exclude $y$, and
replay must reject any escaped local variable.

Constructing a value of \code{Fin n} requires a value \emph{and} a bound proof;
over $\N$, a construction using a signed intermediate representation carries a
nonnegativity obligation. A witness is a program with a proof obligation.

The search grammar should be type-directed and size-bounded: local variables,
constructors, projections, selected recursive calls, affine combinations, and
guarded division or remainder operations. Piecewise terms are permitted only
with exhaustive, proved branch conditions --- never on the assumption that
finitely many sampled regions cover the domain.

\subsubsection*{Affine witnesses with Farkas multipliers}
\status{Implemented and tested in Python.}

Let the input domain be the rational polyhedron $g(x)=Ax+b\ge0$ and the target
constraints $h(x,y)=Cx+Dy+e\ge0$, all componentwise and non-strict, and seek
$y=Mx+c$. A sufficient certificate is a nonnegative matrix $\Lambda$ and a
nonnegative vector $s$ with
\begin{align}
C+DM&=\Lambda A,\label{eq:farkas-linear}\\
e+Dc&=\Lambda b+s.\label{eq:farkas-constant}
\end{align}

\begin{proposition}[Affine witness certificate]
If \eqref{eq:farkas-linear}--\eqref{eq:farkas-constant} hold with
$\Lambda,s\ge0$, then for every $x$ with $Ax+b\ge0$ the term $y=Mx+c$ satisfies
$Cx+Dy+e\ge0$.
\end{proposition}
\begin{proof}
Substitution gives $Cx+D(Mx+c)+e=\Lambda(Ax+b)+s$. Every entry on the right is a
sum of nonnegative multiples of nonnegative quantities plus a nonnegative slack.
\end{proof}

The decisive observation is that \eqref{eq:farkas-linear} and
\eqref{eq:farkas-constant} are \emph{linear in the unknowns} $M,c,\Lambda,s$,
because $A,b,C,D,e$ are fixed problem data. Witness synthesis and proof search
are therefore \emph{one} linear feasibility problem rather than an enumeration
of candidate affine functions followed by a separate verification pass. With $n$
inputs, $m$ outputs, $r$ domain constraints and $q$ target constraints the
encoding has $m(n+1)+q(r+1)$ unknowns and $q(n+1)$ coefficient equations before
template restrictions; entries of $M,c$ are unrestricted and entries of
$\Lambda,s$ are constrained nonnegative. A constant-only grammar adds $M=0$.

For a nonempty rational polyhedron, the affine implication form of Farkas'
lemma gives a completeness route for a \emph{fixed} affine witness: validity of
each non-strict affine target inequality has such
multipliers~\cite{mosek-cookbook}. The floating-point proposer is not a complete
exact LP algorithm and an empty domain needs separate handling, but neither
qualification affects certificate soundness.

A variant fragment solves $A w=Bx+c$ for $w=Wx+d$ by coefficient matching,
giving the exact systems $AW=B$ and $Ad=c$, each solved by rational reduced
row-echelon elimination with free variables set to zero, and re-multiplied
exactly by the checker with no sampled values of $x$ involved.

\subsubsection*{Integer witnesses are a different problem}

A rational affine solution is not an integer proof. The implemented integer mode
merely requires every coefficient of $W$ and $d$ to be integral, and is
deliberately incomplete: a different choice of free variables might yield an
integral solution when the default rational one does not. The instructive case
is $2w=x$, which has the rational witness $x/2$ and no integer witness valid for
every integer $x$; an integer solver must obtain a parity assumption or fail,
and must never silently reinterpret rational division as integer arithmetic.

A practical first integer grammar is
\[
 w(x) ::= c+\sum_i a_ix_i\;\Big|\;
 \left\lfloor\frac{c+\sum_i a_ix_i}{d}\right\rfloor\;\Big|\;
 \operatorname{ite}(g,w_1,w_2),\qquad d>0,
\]
with every candidate checked in the source integer semantics.

\subsubsection*{Integer lattices and certificates of impossibility}
\status{Implemented and tested in Python. The only worker here that emits a
certificate of infeasibility.}

A stronger integer fragment solves $Ax=b$ over $\Z$ by row-by-row unimodular
elimination, returning either a particular solution with a basis for all
homogeneous freedom, or a replayable divisibility obstruction. It does not
search a bounded numerical box. For $6x+10y=2$ it returns
\[
(x,y)=(2,-1)+t(-5,3),\qquad t\in\Z,
\]
and the trace explains why this family contains \emph{all} solutions --- which
is what a later impossibility argument needs.

After processing rows $0,\ldots,i-1$, maintain
\begin{equation}\label{eq:lattice}
\{x\in\Z^n:A_{<i}x=b_{<i}\}=\{x_0+Bz:z\in\Z^d\},
\end{equation}
initially $x_0=0$, $B=I_n$, $d=n$. For the next row $a$ compute $v=aB$ and
$r=b_i-ax_0$, and find an integer unimodular $U$ with
$vU=(g,0,\ldots,0)$, $g\ge0$. If $g=0$, consistency requires $r=0$ and the
family is unchanged. If $g>0$, consistency requires $g\mid r$; with $q=r/g$,
update
\[
x_0\leftarrow x_0+(BU)_{:,0}\,q,\qquad
B\leftarrow(BU)_{:,1:},\qquad d\leftarrow d-1 .
\]
The elementary block, for a current pair $(a,b)$ with $sa+tb=g$, is
\[
\begin{pmatrix}s&-b/g\\t&a/g\end{pmatrix},
\]
whose determinant is $(sa+tb)/g=1$ and which sends the pair to $(g,0)$; a sign
flip makes the final first entry nonnegative.

\begin{proposition}[Lattice elimination]
Each successful row update preserves \eqref{eq:lattice}, and a reported
divisibility obstruction proves that no integer solution exists.
\end{proposition}
\begin{proof}
Because $U$ is unimodular, $z=Uw$ is a bijection of $\Z^d$. The new row equation
becomes $gw_0=r$. For $g>0$ this has an integer solution precisely when
$g\mid r$, in which case $w_0=r/g$ and all other coordinates remain free;
substitution gives the stated update. For $g=0$ the row is either redundant or
contradictory. Induction over processed rows proves both claims.
\end{proof}

The replay format stores, per row, the row number, $U$, $g$, $r$, and either the
quotient or an obstruction marker. The independent checker recomputes the
current row image and residual, checks that all entries are integers, verifies
$\det U=\pm1$, and performs the exact state update --- without calling the
extended-gcd construction or the solver. For $[6\ 10]$ one trace step is
\[
U=\begin{pmatrix}2&-5\\-1&3\end{pmatrix},\qquad
[6\ 10]\,U=[2\ 0],\qquad \det U=1,
\]
and with right-hand side $1$ the same reduction establishes the obstruction
$2\nmid1$.

A production Lean checker might instead accept elementary column operations or
an explicitly checked integer inverse; the elementary-operation format is
attractive because it has short local correctness lemmas and avoids determinant
expansion. For a positive existential goal, checking a returned witness directly
is cheaper than verifying completeness of the whole lattice --- the latter is
needed only when the family or a negative result is used.

Parameterised right-hand sides extend the same machinery: if $Au_0=b_0$ and
$AU=C$, then $x=u_0+Up$ is a witness family for $Ax=b_0+Cp$. This sufficient
method misses witnesses that need piecewise divisibility conditions or integer
division, so it must not report impossibility when the coefficientwise
construction fails. Carrying affine residuals through the same row algorithm
produces obligations $g\mid r(p)$ and a checked integer quotient once they are
proved --- for example, $2x=p$ requires evenness of $p$. If the goal also
includes $Cx\le d$, substitute the lattice parameterisation to obtain
inequalities in the remaining integer parameters; \emph{dropping the
inequalities after solving the equalities would be unsound}, and
natural-number witnesses additionally require nonnegativity, since solutions
over $\Z$ do not automatically inhabit $\N$.

There are at most $n$ dimension-reducing steps, but intermediate coefficient
bit lengths rather than matrix dimensions can dominate cost. The prototype is a
straightforward exact implementation, not an optimised Smith or Hermite
normal-form library.

\subsubsection*{Finite-residue witnesses}
\status{Implemented and tested in Python.}

A different existential shape is discharged by finite case analysis on a residue
class. Given a modulus $m$, synthesise offsets $d_r$, one per residue class
$r$ modulo $m$, such that for every integer $n$,
\[
 n\le n+d_{n\bmod m}<n+m
 \quad\text{and}\quad
 a\,(r+d_r)\equiv b \pmod m .
\]
Verification is a finite check over the $m$ residue classes; the universal
interpretation is justified by the quotient/remainder theorem, not by the
sampled integers. When the congruence $a\,z\equiv b\pmod m$ has no solution ---
detected by a gcd divisibility condition --- the worker reports impossibility for
that parameter tuple rather than failing silently. The implementation is
exercised at inputs such as $10^{20}+7$ to confirm that the finite table really
does instantiate at arbitrary integers.

\subsubsection*{Polynomial witnesses certified by nonnegativity}
\status{Implemented and tested in Python.}

The composition pattern the whole architecture argues for appears in miniature
here. For
\[
 \forall x\in\R,\quad \exists y\in\R,\quad x\le y\ \wedge\ -x\le y,
\]
one prototype enumerates $w(x)=a+bx+cx^2$ with $a,b,c\in\{-2,\ldots,2\}$,
ordered by coefficient $\ell_1$ size and then lexicographically, using five
sample values only as a rejection filter. The 23rd candidate is $w(x)=x^2+1$,
and acceptance requires \emph{two} nonnegativity certificates, for $w-x\ge0$ and
$w+x\ge0$. Convenient explanatory decompositions are
\[
 x^2+1-x=\left(x-\tfrac12\right)^2+\tfrac34,\qquad
 x^2+1+x=\left(x+\tfrac12\right)^2+\tfrac34,
\]
but the recorded certificates contain the different rational weighted-square
decompositions the LP actually returned; both replay exactly. The artefact does
not pretend that a hand-simplified certificate was the optimiser's output.

A production implementation should also try simpler witnesses --- an available
absolute-value or maximum expression --- when the local library makes those
cheap. Restricting a prototype to polynomials is an experiment design, not a
recommended restriction on all \forge{} witnesses.
